\section{The first acceptance cases and the research record}
\label{sec:evidence}

The proposal can be made concrete before connecting to any bank system. The
following cases specify the initial demonstration's intended behavior. Monetary
inputs refer to the adopted financial definitions; the table does not imply that
those facts alone establish SPI status. These are source-based design cases for
independent compliance review, not executed test results.

\begingroup\small
\begin{longtable}{P{0.09\textwidth}P{0.41\textwidth}P{0.4\textwidth}}
\caption{Initial acceptance scenarios for the selected scope.}
\label{tab:cases}\\
\toprule Case & Synthetic fact or source change & Required demonstration \\\midrule
\endfirsthead
\toprule Case & Synthetic fact or source change & Required demonstration \\\midrule
\endhead
F1 & Portfolio HK\$40m; qualifying net assets HK\$0. & Establish the financial
condition via the inclusive portfolio comparison.\\
F2 & Portfolio HK\$35m; qualifying net assets HK\$90m. & Establish the financial
condition via the alternative net-asset route.\\
F3 & Portfolio unknown; qualifying net assets HK\$70m. & Leave the financial
condition unresolved; identify missing portfolio evidence.\\
F4 & Portfolio unknown; qualifying net assets HK\$80m. & Establish the financial
condition via the known sufficient route.\\
F5 & Assets HK\$120m, liabilities HK\$15m, primary residence HK\$30m. & Derive
HK\$75m for the illustrated net-asset calculation; assess the separate portfolio
route independently.\\
O1 & HK\$60m joint portfolio with a non-associate; written client allocation 20\%. &
Use HK\$12m for this portfolio-share illustration; explain the agreement and actor roles.\\
K1 & Transaction-experience route established only for one category; a different
category proposed. & Do not propagate a global experience flag; inspect the
alternative qualification routes and selected categories.\\
C1 & A qualifying wealth record, but conservative investment objectives. & Explain
the investment-objective exclusion; wealth alone does not establish SPI treatment.\\
E1 & Designated-account monitoring selected; exposure rises above threshold without
new funds. & Show the relevant FAQ conditions and restrictions; avoid a universal
transaction-block rule.\\
E2 & A withdrawal instruction precedes a proposed streamlined decision in the
adopted event order. & Use the changed consent state and retain the original
instruction and ordering evidence.\\
T1 & A third-party verification clause with an SFC-request trigger. & Preserve the
trigger; distinguish an absent request from missing request records.\\
T2 & Old tokenisation source replaced by the 2026 circular. & Preserve both editions,
record replacement, identify affected rules and reproduce historical decisions.\\
M1 & Product-linked gift versus a discount of fees or charges. & Preserve the stated
exception and the scope of the cited marketing provision.\\
M2 & Daily dealing fund with weekly redemption versus a different daily cut-off. &
Distinguish reduced dealing frequency from administrative procedure.\\
D1 & A thematic-review announcement enters the source collection. & Classify and
track the announcement without inventing a new numerical eligibility criterion.\\
\bottomrule
\end{longtable}
\endgroup

Cases F1--F5 and O1 derive from Annex 1, paragraphs 3.1--3.4. K1 and C1 concern
paragraphs 4--7. E1 requires paragraph 8.3 together with Annex 2, FAQ 6; E2 concerns
the consent provisions and the explicitly adopted event-ordering model. T1--T2
refer to the cited tokenisation editions; M1--M2 to the marketing circular; D1
to the thematic-review announcement
\citep{sfcspiannex1,sfcspiannex2,sfctoken2023,sfctoken2026,sfcmarketing,sfcreview2025}.
This source mapping is part of the acceptance package and should be visible
beside each case in the workbench.

\subsection{Conformance cases and the research record}

The proposal now has a machine-readable companion to the business scenarios.
Under \path{docs/specs/v0.1/fixtures/}, 35 decision cases specify complete bundles,
synthetic snapshots, assessment times and expected result projections. Six
deliberately invalid variants exercise schema, reference, type and identifier
failures. Eight event histories cover withdrawal, regrant, timely performance,
deadline equality, missing observations and late performance. Four release cases
distinguish valid approval from self-review, open issues and stale evidence.

These sets serve different purposes. The 35 cases specify the selected language;
the fifteen business scenarios above identify regulatory interpretation and
operational questions, several of which remain beyond the first executable
slice. An agent must not claim that passing the language cases completes those
business scenarios. For example, a Boolean assessment input can represent a
reviewer's judgment about investment objectives, but it does not automate that
judgment. Likewise, ownership and category definitions require their own reviewed
normalization and end-to-end tests before release.

The delivered checker validates all three JSON schemas, the 35 decision fixture
contracts, six negative examples, event-record structure, original/derivative/span
hashes and the SQLite table definitions. It does not execute the planned decision,
event or release engines. Its report says so explicitly. W03, W04 and W05 provide
the corresponding runtime acceptance commands. Independent legal adjudication,
human usability assessment and published-tool proof replay remain prospective.

\subsubsection{A reproducible research collection}

The papers are stored under \path{docs/papers/} using the requested filename
convention: the title, a comma, the first author's surname and the year in
parentheses. Filesystem-unsafe title punctuation is normalized, while the manifest
retains the original title. Each
record includes its source URL, bibliographic identity, edition note, retrieval
date, byte count, PDF page count and SHA-256 digest. The library index distinguishes
distinct works from alternate editions and identifies discovery leads that were
not available as usable full text.

The technical notes identify which part of each paper supports a proposed
borrowing and what remains untested. They distinguish a semantics from its
implementation, a reported experiment from a reproduced result, and a publicly
available repository from a licensed and approved bank dependency. Full-text
extraction is retained for local search, but the source PDF remains the reference
where layout or equations matter. ResearchAssistant extraction is an aid to
inspection, not certification of the extracted text.

The SFC source snapshots and the earlier local-repository inspection records
remain under \path{.localresources/}, with their own manifest. The LaTeX and
bibliography are in \path{docs/proposal/}. This arrangement permits another
reviewer to trace the proposal's substantive examples back to a document edition
without reconstructing the web searches. Redistributing the paper library is a
separate matter from using the local copies for this research; the source licenses
and access terms remain relevant.

\subsection{Decisions needed at the start of a bank prototype}

The first decision is the precise business perimeter: legal entity, activity,
client class and product family. The second is the owner of each interpretation
and the independent review process for disputed cases. The third is the mapping
from regulatory definitions to actual bank records, including freshness and
missing-evidence treatment. These decisions should be made on the selected source
slice, with worked examples, rather than as abstract architecture questions.

The initial technical commitment should be modest and testable: one source
repository, one typed decision semantics, one bounded event model, one review
interface and one deployment adapter. The literature gives substantial guidance
for each part. It also explains why a universal translation-and-proof promise
would be unsound. A successful prototype would show which interpretations are
adopted, which computations follow them, which facts support the result, and where
human judgment is still required.
